\documentclass[11pt]{article}
\usepackage{setspace}
\usepackage{geometry}
\geometry{margin=1in}
\setstretch{1.15}

\begin{document}

\noindent
\textbf{Zhongli Wang} \\
Phone: +49 160 9541 6015 \\
Email: nkwangzhongli@gmail.com \\
Website: https://sites.google.com/view/zhongli-wang \\
Bielefeld, Germany \\

\vspace{0.5cm}

\noindent
\textbf{Prof.\ Ferdinand Vieider} \\
RISLab, Department of Economics \\
Ghent University \\
Ghent, Belgium \\

\vspace{0.5cm}

\noindent
\textbf{Postdoc ``Noisy Cognition as a Unified Driver of Decisions''} \\

\vspace{0.3cm}

\noindent Dear Prof.\ Vieider,\\[0.2cm]

I am writing to apply for the postdoctoral position in RISLab on the project \textit{Noisy Cognition as a Unified Driver of Decisions}. I am currently completing my Ph.D.\ in Economics at Bielefeld University and Università Cattolica del Sacro Cuore within the Marie Sk\l{}odowska--Curie EPOC programme. My research is quantitative and behavioral in spirit: I study how bounded rationality, belief updating, and heterogeneous heuristics generate systematic deviations from benchmark predictions and how these deviations aggregate in dynamic environments. The project’s goal---to develop stochastic, Bayesian models that unify decision patterns across domains (risk, uncertainty, time; gains, losses, mixed gambles) and to test them in large-scale empirical work---matches both my methodological strengths and my research interests.

My doctoral training emphasized formal modeling and computational inference in settings where decision rules are noisy and information is imperfect. In my dissertation, I build dynamic models with heterogeneous agents who learn from local information, update beliefs, and adapt strategies under constraints. Although my main empirical setting is technology-intensive markets, the underlying mechanisms---noisy perception, biased inference, and heuristic choice---are directly aligned with your project’s focus on cognitive drivers of decisions across contexts. I am comfortable translating verbal behavioral hypotheses into precise stochastic structures, stress-testing identification with simulation, and producing transparent, reproducible code.

I also have early experience with behavioral research questions and individual-level choice analysis. During my Master’s studies at the University of Zurich, I completed coursework in Experimental Economics and wrote a semester paper on the link between (over/under)confidence and procrastination, highlighting how biased self-beliefs correlate with dynamically inconsistent behavior. This work sharpened my interest in belief distortions and the cognitive origins of seemingly disparate behaviors---a perspective that closely fits a unifying ``noisy cognition'' approach.

Methodologically, I can contribute to the project’s core deliverables in three ways. First, I can help develop and implement Bayesian stochastic choice models that discipline noise in perceptions, preferences, and memory, and that generate comparable primitives across domains. Second, I can support the design and analysis of large-scale empirical studies by building end-to-end pipelines for experimental data collection, cleaning, preregistered analysis, and model-based estimation. Third, I can contribute computational expertise for inference and validation (e.g., hierarchical modeling, posterior predictive checks, simulation-based calibration, and sensitivity analyses), using Python and/or R for statistical coding and reproducible workflows. While I am not a Bayesian specialist by field identity, I have strong quantitative foundations and substantial experience with simulation-based estimation, which positions me to become highly productive with Bayesian techniques and to apply them rigorously in this setting.

Beyond my dissertation, I have coauthored work that demonstrates my ability to build robust quantitative pipelines and handle high-dimensional data. \textit{Polytope Fraud Theory} (\textit{International Review of Financial Analysis}, 2025) applies supervised learning to detect structured irregularities in firm financial statements. A second project, \textit{Agentic Workflows for Economic Research}, develops reproducible research pipelines that integrate large-language-model tools with careful documentation and evaluation. While these projects address different substantive questions, they reflect the same strengths that matter for this postdoc: statistical coding discipline, careful model evaluation, and the ability to connect quantitative outputs to interpretable mechanisms.

If selected, I would be excited to pursue a complementary independent research direction alongside the project, focused on connecting noisy cognition to technology-mediated decision environments. One example is to study whether interacting with algorithmic decision aids (e.g., recommender systems or AI assistants) attenuates or amplifies cognitive noise, and whether this differs systematically across domains (risk vs.\ time vs.\ ambiguity) or across populations. This would be a natural extension of the project’s unifying perspective and could generate additional testable predictions and publishable results.

Regarding eligibility, I understand that this fellowship is awarded under international mobility conditions. I have been based in Germany and Italy during the last years of my PhD and have not worked or studied in Belgium during the relevant period, aside from short visits such as conferences. I would be happy to provide any additional documentation needed to confirm compliance with the mobility rule and to discuss a feasible start date (flexible) depending on project timing and my PhD completion.

Thank you for considering my application. I would welcome the opportunity to discuss how my quantitative background, modeling experience, and interest in belief biases can contribute to the goals of \textit{Noisy Cognition as a Unified Driver of Decisions} and to the research environment at RISLab.\\[3mm]

\noindent Sincerely,\\[12mm]
\noindent Zhongli Wang

\vspace{4mm}

% \noindent \textit{Potential referees (available upon request):} \\
% \noindent Prof.\ Herbert Dawid (Bielefeld University) \\
% \noindent Prof.\ Domenico Delli Gatti (Università Cattolica del Sacro Cuore)

\end{document}